\section{Common moment estimates and the SCAFFOLD proof}
\label{app:scaf-proof}

The proof separates the moment estimates shared by both algorithms from the
additional contribution of SCAFFOLD controls. We first bound local fluctuations,
then obtain the sharp global raw second moment and a stationary mean-balance identity
with a uniform remainder. The SCAFFOLD control covariance completes
Theorem~\ref{thm:scaf-main}; Appendix~\ref{app:fa-proof} uses the same
estimates for FedAvg and the matched attribution.

\subsection{Notation and stationary inputs}
\label{app:proof-notation}
\label{app:common-inputs}

The three scales used to track errors uniformly in $N$ are
\begin{equation}
\label{eq:scales}
 s:=s_{\gamma,N}=\frac{\gamma}{N}+\gamma^2,\qquad
 r:=r_{\gamma,N}=\frac{\gamma^2}{N}+\gamma^3=\gamma s,\qquad
 v:=v_{\gamma,N}=\frac{\gamma^2}{N^2}+\gamma^3.
\end{equation}
Constants denoted by $K$ may change between occurrences and depend on fixed
$H,d$ and the objective/noise regularity constants, but not on $N$ or $\gamma$.
All $O(\cdot)$ terms in these proofs have this same uniform meaning.
We use Euclidean norms for vectors and tensors, induced norms for linear
maps, and $\|\cdot\|_F$ for the Frobenius matrix norm. In particular,
$\|P^h\|$ is an operator norm and operators on matrices use the
Frobenius-induced norm. We use repeatedly, for
$\gamma\le1$ and $N\ge2$,
\begin{equation}
\label{eq:common-scale-comparisons}
 v\le r,\qquad \gamma^2\sqrt{s}\le r,\qquad
 \gamma^2(s+N^{-1})\le2r,\qquad \gamma^4\le\gamma v.
\end{equation}
For example, $\sqrt{s}\le\gamma+N^{-1}$ follows by squaring both sides.
For vectors $u,v$, set $(T[u,v])_i=\sum_{j,k}T_{ijk}u_jv_k$.
In cubic identities, $\Sym$ is the unnormalized sum of three placements:
\[
 [\Sym(u\otimes M)]_{ijk}=u_iM_{jk}+u_jM_{ik}+u_kM_{ij},\qquad
 [\Sym(u^{\otimes2}\otimes v)]_{ijk}
 =u_iu_jv_k+u_iv_ju_k+v_iu_ju_k,
\]
with $M$ symmetric. This convention extends linearly under expectation.

We use stationary moment bounds from \citet[Corollary~B.4 and
Lemma~B.5]{mangold2025scaffold} and \citet[Proposition~3 and
Lemmas~5--6]{mangold2025fedavg}. These results apply here because the
noise laws are fixed across rounds, $F_\eps(y)=f(y)+\langle\eps,y\rangle$
has the same Hessian as $f$, and bounded additive noise has the required
moments without state-dependent growth. Identical client objectives remove
the heterogeneity term. A common sufficient step-size threshold is
\begin{equation}
\label{eq:imported-step-threshold}
 0<\gamma\le\gamma_{\mathrm{imp}}
 :=\min\left\{1,\frac1{48L},\frac1{24H(L+\mu)}\right\}.
\end{equation}
It also implies the FedAvg higher-moment restrictions $\gamma\le1/(9L)$
and $\gamma\mu H\le1$. The additional SCAFFOLD variance-growth restrictions
are vacuous because their parameter $\beta$ is zero here. We decrease
$\gamma_0\le\gamma_{\mathrm{imp}}$, independently of $N$, for subsequent
contraction and absorption inequalities.

For each algorithm under its own stationary law, the imported estimates give
\begin{equation}
\label{eq:A-coarse-even}
 \E\|x\|^{2p}+\frac1N\sum_c\E\|\theta_{c,h}\|^{2p}
 =O(\gamma^p),\qquad p=1,2,3,\quad 0\le h\le H.
\end{equation}
For SCAFFOLD, additionally, we have
\begin{equation}
\label{eq:A-control-sixth}
 \sup_c\E\|\xi_c\|^6\le K,\qquad \sum_c\xi_c=0\quad\text{pathwise}.
\end{equation}
In the common lemmas below, write $\xi_c=0$ for FedAvg. This notation
unifies only the within-round local recursions; it does not impose the
SCAFFOLD control-refresh rule on FedAvg. In either case, controls are fixed
during a round, and fresh noise is independent of the complete round-start state.

We set $\delta_{c,h}=\theta_{c,h}-x$, $\bar\delta_h=N^{-1}\sum_c\delta_{c,h}$,
$\bar\eps_j=N^{-1}\sum_c\eps_{c,j}$, and
\begin{equation}
\label{eq:A-qP}
 \begin{aligned}
 q(y)&=\nabla f(y)-Ay,\qquad P=I-\gamma A,\qquad b=\E[x],\\
 m_h&=\frac1N\sum_c\E[\theta_{c,h}],\qquad
 \nu_h=\frac1N\sum_c\E[q(\theta_{c,h})].
 \end{aligned}
\end{equation}
Also, we write
\begin{equation}
\label{eq:A-MU}
 M_x=\E[xx^\top],\qquad M_h=\frac1N\sum_c\E[\theta_{c,h}\theta_{c,h}^\top],
 \qquad U_h=M_h-M_x,\qquad Q_\xi=\frac1N\sum_c\E[\xi_c\xi_c^\top].
\end{equation}
These are raw second moments, evaluated under each algorithm's own
stationary law. In particular, $U_h$ measures the additional second moment
accumulated during the $h$ local steps.

\subsection{Common coarse and local second-moment estimates}
\label{app:common-coarse}

\begin{lemma}[Coarse scales and local second moments]
\label{lem:common-coarse}
For either algorithm and $0\le h\le H$,
\begin{align}
 \frac1N\sum_c\E\|\delta_{c,h}\|^{2p}&=O(\gamma^{2p}),\quad p=1,2,3,
 \label{eq:A-delta-moments}\\
 \E\|x\|^2+\|b\|&=O(s),\qquad m_h-b=O(r),
 \label{eq:A-mean-bootstrap}\\
 \E\|\bar\delta_h\|^2&=O\bigl(\gamma^2(s+N^{-1})\bigr),
 \label{eq:A-deltabar}\\
 U_h&=\gamma^2(hC+h^2Q_\xi)+O(r).
 \label{eq:A-U-pre}
\end{align}
\end{lemma}

\begin{proof}
Unrolling the fixed-$H$ local recursion and using Equations~\eqref{eq:A-coarse-even}
and~\eqref{eq:A-control-sixth} gives Equation~\eqref{eq:A-delta-moments}
(the control bound is automatic for FedAvg).
At the server, the client control corrections cancel by the zero-sum
condition. The resulting round increment and stationary balance are
\begin{align}
 x^+-x&=-\gamma\sum_{h=0}^{H-1}\frac1N\sum_c\nabla f(\theta_{c,h})
         -\gamma\sum_{j=1}^H\bar\eps_j,
 \label{eq:A-global-increment}\\
 0&=\sum_{h=0}^{H-1}\frac1N\sum_c\E[\nabla f(\theta_{c,h})].
 \label{eq:A-balance}
\end{align}
\emph{Coarse global moment and mean.}
Write $x^+=\Phi(x)+\zeta+E_{\mathrm{traj}}$, where
$\Phi(x)=x-\gamma H\nabla f(x)$, $\zeta=-\gamma\sum_{j=1}^H\bar\eps_j$,
and $E_{\mathrm{traj}}=-\gamma\sum_hN^{-1}\sum_c
\{\nabla f(\theta_{c,h})-\nabla f(x)\}$.
Strong convexity gives $\|\Phi(x)\|^2\le(1-c\gamma)\|x\|^2$ for some $c>0$,
while $\E\|\zeta\|^2=O(\gamma^2/N)$ and
$\E\|E_{\mathrm{traj}}\|^2=O(\gamma^4)$.
First, we use $\E\langle\Phi(x),\zeta\rangle=0$, then apply Young's inequality
only to $E_{\mathrm{traj}}$. For fixed sufficiently small $a>0$, we derive
\[
 \E\|x\|^2\le(1+a\gamma)
 \{(1-c\gamma)\E\|x\|^2+K\gamma^2/N\}
 +K_a\gamma^{-1}\E\|E_{\mathrm{traj}}\|^2.
\]
Absorption gives $\E\|x\|^2=O(s)$ and hence
$N^{-1}\sum_c\E\|\theta_{c,h}\|^2=O(s)$.
Since $\|q(y)\|\le K\|y\|^2$, we have $\nu_h=O(s)$.
The averaged recursion $m_{h+1}=Pm_h-\gamma\nu_h$, with $m_0=m_H=b$, gives
\[
 (I-P^H)b=-\gamma\sum_{j=0}^{H-1}P^{H-1-j}\nu_j,\qquad
 m_h-b=(P^h-I)b-\gamma\sum_{j=0}^{h-1}P^{h-1-j}\nu_j.
\]
Using $\|(I-P^H)^{-1}\|=O(\gamma^{-1})$ and $\|P^h-I\|=O(\gamma)$
proves Equation~\eqref{eq:A-mean-bootstrap}. For the client-averaged displacement, the controls again cancel. Its
gradient part has second moment
$O(\gamma^2s)$ and the averaged-noise part has second moment $O(\gamma^2/N)$.
Their sum gives Equation~\eqref{eq:A-deltabar} without assuming the two parts
are independent.

\emph{Local second moments.}
To isolate the effect of fresh noise and persistent controls, write
$\theta_{c,h}=x+D_{c,h}+R_{c,h}$ with
\begin{equation}
\label{eq:A-DR}
 D_{c,h}=-\gamma\left(h\xi_c+\sum_{j=1}^h\eps_{c,j}\right),\qquad
 R_{c,h}=-\gamma\sum_{\ell=0}^{h-1}\nabla f(\theta_{c,\ell}).
\end{equation}
Zero-sum cancellation and freshness imply
\begin{equation}
\label{eq:A-DD}
 \frac1N\sum_c\E[xD_{c,h}^\top]=0,\qquad
 \frac1N\sum_c\E[D_{c,h}D_{c,h}^\top]=\gamma^2(hC+h^2Q_\xi).
\end{equation}
Only the client-averaged first identity is needed; $x$ and individual controls
need not be uncorrelated. Furthermore,
$N^{-1}\sum_c\E\|D_{c,h}\|^2=O(\gamma^2)$ and
$N^{-1}\sum_c\E\|R_{c,h}\|^2=O(\gamma^2s)$.
Cauchy--Schwarz bounds the averaged $xR^\top$, $DR^\top$, and $RR^\top$
terms by $O(\gamma s)$, $O(\gamma^2\sqrt{s})$, and $O(\gamma^2s)$,
respectively. All are $O(r)$ by Equation~\eqref{eq:common-scale-comparisons},
which proves Equation~\eqref{eq:A-U-pre}.
\end{proof}

\subsection{Common higher-moment bootstrap}
\label{app:common-higher}

The signed third moment is needed to control the cubic term in the
stationary Taylor expansion. We obtain it after a fourth-moment estimate.

\begin{lemma}[Fourth and signed third moments]
\label{lem:common-higher}
For either algorithm,
\begin{equation}
\label{eq:A-higher-moments}
 \E\|x\|^4=O(v),\qquad \|\E[x^{\otimes3}]\|=O(v).
\end{equation}
\end{lemma}

\begin{proof}
Unrolling the averaged local recursion gives
\begin{equation}
\label{eq:A-round-exact}
 x^+=P^Hx+\eta+\Delta,\qquad
 \eta=-\gamma\sum_{j=1}^HP^{H-j}\bar\eps_j,\quad
 \Delta=-\gamma\sum_{h=0}^{H-1}P^{H-1-h}\bar q_h,
 \quad \bar q_h=\frac1N\sum_c q(\theta_{c,h}).
\end{equation}
The centered fresh term $\eta$ is independent of the round-start state.
Expansion of sums of independent centered vector blocks gives
\begin{equation}
\label{eq:A-eta-moments}
 \E\|\eta\|^2=O(\gamma^2/N),\qquad
 \E\|\eta\|^4=O(\gamma^4/N^2),\qquad
 \|\E[\eta^{\otimes3}]\|=O(\gamma^3/N^2).
\end{equation}
No symmetry of the noise law is used. Global Lipschitzness of $\nabla f$
and bounded $\nabla^3f$ imply $\|q(y)\|\le K\min\{\|y\|,\|y\|^2\}$.
Thus $\|q(y)\|^2\le K\|y\|^4$ and $\|q(y)\|^4\le K\|y\|^6$;
Jensen's inequality and Equation~\eqref{eq:A-coarse-even} yield
\begin{equation}
\label{eq:A-Delta-moments}
 \E\|\Delta\|^2=O(\gamma^4),\qquad \E\|\Delta\|^4=O(\gamma^7).
\end{equation}
In particular, an eighth-moment assumption is unnecessary.

\emph{Fourth moment.}
Let $Y=P^Hx+\eta$ and $M_4=\E\|x\|^4$.
Independence and centeredness of $\eta$ give
\[
 \E\|Y\|^4\le\|P^H\|^4M_4+
 K\left\{\E\|x\|^2\E\|\eta\|^2
 +\|b\|\sqrt{\E\|\eta\|^2\E\|\eta\|^4}+\E\|\eta\|^4\right\}
 \le\|P^H\|^4M_4+K\gamma v.
\]
Use $\|y+z\|^4\le(1+a\gamma)\|y\|^4+K_a\gamma^{-3}\|z\|^4$
and $\|P^H\|^4\le1-c\gamma$. For fixed sufficiently small $a>0$,
stationarity implies
\[
 M_4\le(1+a\gamma)\{(1-c\gamma)M_4+K\gamma v\}
      +K_a\gamma^{-3}\E\|\Delta\|^4.
\]
The last term is $O(\gamma^4)\subset O(\gamma v)$; absorption proves
$M_4=O(v)$ and also $\E\|Y\|^4=O(v)$.

\emph{Signed third moment.}
For $T_3=\E[x^{\otimes3}]$, the linear part satisfies
\[
 \E[Y^{\otimes3}]=(P^H)^{\otimes3}T_3
 +\Sym\{P^Hb\otimes\E[\eta\eta^\top]\}+\E[\eta^{\otimes3}].
\]
The two forcing terms are $O(\gamma v)$ by the coarse mean bound and
Equation~\eqref{eq:A-eta-moments}.
Set $a_h=N^{-1}\sum_c\|\theta_{c,h}\|^2$. By the fourth-moment bound
just proved and Equation~\eqref{eq:A-delta-moments},
\[
 \E[a_h^2]\le\frac1N\sum_c\E\|\theta_{c,h}\|^4
 \le K(v+\gamma^4)=O(v),\qquad
 \E[\|Y\|^2a_h]\le\sqrt{\E\|Y\|^4\E[a_h^2]}=O(v).
\]
Since $\|\Delta\|\le K\gamma\sum_ha_h$, it follows that
$\|\E[Y^{\otimes2}\otimes\Delta]\|=O(\gamma v)$.
The other terms in the cubic expansion obey
\[
 \E[\|Y\|\|\Delta\|^2]\le
 \sqrt{\E\|Y\|^2\E\|\Delta\|^4}=O(\gamma^4),\qquad
 \E\|\Delta\|^3\le(\E\|\Delta\|^4)^{3/4}=O(\gamma^{21/4}),
\]
both $O(\gamma v)$. Hence,
$(I-(P^H)^{\otimes3})T_3=O(\gamma v)$; its inverse has norm
$O(\gamma^{-1})$, proving $T_3=O(v)$. The sharp global second moment can
therefore be established using these higher-moment bounds.
\end{proof}

\subsection{Common sharp global second moment}
\label{app:common-covariance}

We now compare the nonlinear stationary chain with its linear counterpart.
The leading global second moment comes from server-averaged noise; the
higher-moment bounds control the nonlinear correction.

\begin{lemma}[Sharp stationary raw second moment]
\label{lem:A-sharp-second-moment}
For either algorithm, the nonlinear forcing in the stationary second-moment
equation is $O(\gamma r)$, and
\begin{equation}
\label{eq:A-sharp-Mx}
 M_x=\frac{\gamma}{N}S_C+O(r).
\end{equation}
\end{lemma}

\begin{proof}
\emph{Linear chain.}
For the linear chain $x_{\rm lin}^+=P^Hx_{\rm lin}+\eta$,
$\E[\eta\eta^\top]=(\gamma^2/N)\sum_{k=0}^{H-1}P^kC(P^k)^\top$.
Its stationary second-moment equation is the $H$-fold iteration of
$M_{\rm lin}=PM_{\rm lin}P^\top+\gamma^2C/N$.
Both equations have unique solutions because $P$ is contractive.
In an eigenbasis of $A$, with eigenvalues $a_i$,
\begin{equation}
\label{eq:A-Mlin-coord}
 (M_{\rm lin})_{ij}=\frac{\gamma}{N}
 \frac{C_{ij}}{a_i+a_j-\gamma a_ia_j},\qquad
 (S_C)_{ij}=\frac{C_{ij}}{a_i+a_j}.
\end{equation}
Consequently $M_{\rm lin}=(\gamma/N)S_C+O(\gamma^2/N)$ without
assuming $AC=CA$.

\emph{Nonlinear state correction.}
To control $\E[x\Delta^\top]$, expand
$\bar q_h=q(x)+Dq(x)\bar\delta_h+\bar\rho_h$,
where $\|Dq(x)\|\le K\|x\|$ and
$\|\bar\rho_h\|\le K N^{-1}\sum_c\|\delta_{c,h}\|^2$.
The expansion $q(x)=\frac12T[x,x]+O(\|x\|^3)$ and the preceding lemmas give
\begin{align*}
 \|\E[xq(x)^\top]\|_F&\le K\{\|\E[x^{\otimes3}]\|+\E\|x\|^4\}=O(r),\\
 \E[\|x\|^2\|\bar\delta_h\|]&\le
 \sqrt{\E\|x\|^4\E\|\bar\delta_h\|^2}
 =O\bigl(\sqrt{v\gamma^2(s+N^{-1})}\bigr)=O(r),\\
 \E[\|x\|\|\bar\rho_h\|]&=O(\gamma^2\sqrt{s})=O(r).
\end{align*}
The outer factor $\gamma$ in $\Delta$ therefore yields
\begin{equation}
\label{eq:A-xDelta}
 \|\E[x\Delta^\top]\|_F=O(\gamma r).
\end{equation}

\emph{Fresh-noise/nonlinear correlation.}
To control $\E[\eta\Delta^\top]$, hold the complete round-start state fixed.
Index the $NH$ fresh blocks by $i=(c,j)$ and write
$\eta=\sum_iB_i\eps_i$, with $\|B_i\|\le K\gamma/N$.
Let $\Delta^{(i,0)}$ replace block $\eps_i$ by zero.
Only client $c$ changes before aggregation. Lipschitz propagation gives a
path change of at most $K\gamma\|\eps_i\|$ and hence
\[
 \|\Delta-\Delta^{(i,0)}\|\le K\frac{\gamma^2}{N}\|\eps_i\|,\qquad
 \E[\eps_i(\Delta^{(i,0)})^\top\mid\text{round-start state}]=0.
\]
Each term $\E[B_i\eps_i(\Delta-\Delta^{(i,0)})^\top]$ is
$O(\gamma^3/N^2)$; summing the $NH$ terms gives
\begin{equation}
\label{eq:A-etaDelta}
 \|\E[\eta\Delta^\top]\|_F=O(\gamma^3/N).
\end{equation}
Replacing one noise block at a time handles the dependence between
$\eta$ and $\Delta$. Finally,
$\|\E[\Delta\Delta^\top]\|_F=O(\gamma^4)$ by
Equation~\eqref{eq:A-Delta-moments}.

\emph{Stationary covariance equation.}
For $\mathcal S_\gamma(X)=X-P^HX(P^H)^\top$, stationarity gives
\begin{align}
\label{eq:A-second-moment-difference}
 \mathcal S_\gamma(M_x-M_{\rm lin})
 ={}&P^H\E[x\Delta^\top]+\E[\Delta x^\top](P^H)^\top\notag\\
 &+\E[\eta\Delta^\top]+\E[\Delta\eta^\top]+\E[\Delta\Delta^\top].
\end{align}
The right-hand side is $O(\gamma r)$, since both $\gamma^3/N$ and
$\gamma^4$ are bounded by $\gamma r$. Moreover,
$\|\mathcal S_\gamma^{-1}\|\le(1-\|P^H\|^2)^{-1}=O(\gamma^{-1})$.
Thus $M_x-M_{\rm lin}=O(r)$, proving Equation~\eqref{eq:A-sharp-Mx}.
\end{proof}

\subsection{Common signed Taylor closure}
\label{app:common-taylor}

The global raw second moment determines one part of the mean bias. Stationary
gradient balance also involves the second moments along the local
trajectories. The following lemma combines both through a Taylor expansion.

\begin{lemma}[Local Taylor tail and stationary mean balance]
\label{lem:common-taylor}
For either algorithm,
\begin{align}
 \left\|\frac1N\sum_c\E[\theta_{c,h}^{\otimes3}]\right\|&=O(r),\qquad
 \frac1N\sum_c\E\|\theta_{c,h}\|^4=O(v),
 \label{eq:A-local-moments}\\
 b&=-\frac12A^{-1}T[M_x]
     -\frac1{2H}A^{-1}T\!\left[\sum_{h=0}^{H-1}U_h\right]+O(r).
 \label{eq:A-master}
\end{align}
\end{lemma}

\begin{proof}
Expanding $\theta_{c,h}=x+\delta_{c,h}$ gives
\begin{align*}
 \frac1N\sum_c\E[\theta_{c,h}^{\otimes3}]
 ={}&\E[x^{\otimes3}]+\Sym\{\E[x^{\otimes2}\otimes\bar\delta_h]\}\\
 &+\Sym\!\left\{\frac1N\sum_c\E[x\otimes\delta_{c,h}^{\otimes2}]\right\}
 +\frac1N\sum_c\E[\delta_{c,h}^{\otimes3}].
\end{align*}
The four terms have norms bounded respectively by
\[
 O(v),\qquad O\bigl(\sqrt{v\gamma^2(s+N^{-1})}\bigr),\qquad
 O(\gamma^2\sqrt{s}),\qquad O(\gamma^3),
\]
using Cauchy--Schwarz and Lemmas~\ref{lem:common-coarse}
and~\ref{lem:common-higher}. All are $O(r)$.
Also $N^{-1}\sum_c\E\|\theta_{c,h}\|^4
\le K\{\E\|x\|^4+N^{-1}\sum_c\E\|\delta_{c,h}\|^4\}
=O(v+\gamma^4)=O(v)$.
The signed cubic and absolute fourth-moment bounds now control the
higher-order terms in the gradient expansion. With
$D_4=\nabla^4f(0)$, Taylor's theorem gives
\begin{equation}
\label{eq:A-Taylor}
 \nabla f(u)=Au+\frac12T[u,u]+\frac16D_4[u,u,u]+\rho_5(u),\qquad
 \|\rho_5(u)\|\le K\|u\|^4.
\end{equation}
Here $D_4[u,u,u]$ contracts the fourth derivative against three vectors.
Equation~\eqref{eq:A-local-moments} controls the signed cubic term and
absolute fourth-order remainder, giving
\begin{equation}
\label{eq:A-Taylor-tail}
 \left\|\sum_{h=0}^{H-1}\frac1N\sum_c
 \E\!\left[\nabla f(\theta_{c,h})-A\theta_{c,h}
             -\frac12T[\theta_{c,h},\theta_{c,h}]\right]\right\|=O(r).
\end{equation}
Substitute into Equation~\eqref{eq:A-balance}, use
$\sum_hm_h=Hb+O(r)$ from Lemma~\ref{lem:common-coarse}, and write
$M_h=M_x+U_h$. This gives
$0=HAb+\frac12T[\sum_hM_h]+O(r)$ and proves
Equation~\eqref{eq:A-master}.
\end{proof}

\subsection{SCAFFOLD control moment and coefficient assembly}
\label{app:scaf-assembly}

The common lemmas reduce the algorithm-specific calculation to the
stationary covariance $Q_\xi$ of the persistent controls. Define
\begin{equation}
\label{eq:A-GE}
 G_c=\frac1H\sum_{h=0}^{H-1}
 \left\{\nabla f(\theta_{c,h})-\frac1N\sum_k\nabla f(\theta_{k,h})\right\},
 \qquad E_c=\frac1H\sum_{j=1}^H(\eps_{c,j}-\bar\eps_j).
\end{equation}
The control update in Equation~\eqref{eq:scaf-server} gives
$\xi_c^+=-G_c-E_c$ exactly. Independence across clients and local steps yields
$\E[E_cE_c^\top]=(1-1/N)C/H$. Smoothness and
Equation~\eqref{eq:A-delta-moments} give
$N^{-1}\sum_c\E\|G_c\|^2=O(\gamma^2)$, while
$N^{-1}\sum_c\E\|E_c\|^2=O(1)$.
The $GE^\top$ and $EG^\top$ terms are therefore $O(\gamma)$ by
Cauchy--Schwarz, not by independence. Control stationarity implies
\begin{equation}
\label{eq:A-Q}
 Q_\xi=Q_{\gamma,N,H}=\frac1H\left(1-\frac1N\right)C+O(\gamma).
\end{equation}
Inserting this covariance into the local second-moment expansion adds
the $h^2$ contribution specific to persistent controls:
\begin{equation}
\label{eq:A-U}
 U_h=\gamma^2\left(h+\frac{h^2}{H}\right)C+O(r),
\end{equation}
because the omitted pieces are $O(\gamma^2/N+\gamma^3)$.
Finally,
\begin{equation}
\label{eq:A-U-sum}
 \sum_{h=0}^{H-1}U_h
 =\gamma^2\frac{(H-1)(5H-1)}6C+O(r).
\end{equation}
Combining Equations~\eqref{eq:A-sharp-Mx}, \eqref{eq:A-master},
and~\eqref{eq:A-U-sum} yields
\[
 b^{\Scaf}_{\gamma,N,H}
 =-\frac{\gamma}{2N}A^{-1}T[S_C]
 -\gamma^2\frac{(H-1)(5H-1)}{12H}A^{-1}T[C]+O(r).
\]
This proves Theorem~\ref{thm:scaf-main}, including the stated uniform remainder.


\subsection{Proof of the affine-control translation}
\label{app:affine-proof}
\begin{proof}[Proof of Proposition~\ref{prop:affine-shift}]
The gradient is $\nabla f_c(y)=\nabla f(y)+a_c$. With the translation in Equation~\eqref{eq:affine-objectives}, the heterogeneous local step becomes
\[
\theta_{c,h+1}=\theta_{c,h}-\gamma\{\nabla f(\theta_{c,h})+\widetilde\xi_c+\eps_{c,h+1}\}.
\]
The zero-sum constraint is preserved because $\sum_ca_c=0$, and the refreshed state satisfies
\[
\widetilde\xi_c^+=\xi_c^++a_c
=\widetilde\xi_c+\frac{\theta_{c,H}-x^+}{\gamma H}.
\]
Thus the translated controls give exactly the homogeneous round map,
with the global model $x$ unchanged. Induction yields pathwise equality
for corresponding initial states driven by the same noise blocks.

The translation is invertible and carries invariant laws between the
two chains. Uniqueness of the homogeneous invariant law establishes
equality of their global stationary marginals. Since
$N^{-1}\sum_cf_c=f$, both biases are measured from the same minimizer,
and Theorem~\ref{thm:scaf-main} gives the expansion and uniform remainder.
\end{proof}
